\documentclass{article}
\usepackage{Self_Style}

\title{UChicago REU Notes Week 3}
\author{Zih-Yu Hsieh}
\date{\today}

\begin{document}
\maketitle
\tableofcontents

\section{Operad: Alternative Definition and Relation to Monad (ND, look at algebras)}

For the last section, we've seen an associative non-$\Sigma$ operad, to produce the version in operad (equipped with symmetric group actions), we need to consider another version. But for this, I think there are some technicalities in the category theory, so we'll assume $\mathcal{V}$ is monoidally cocomplete.

An alternative assumption can be made that internal hom functor exists. Out of curiosity, are the two equivalent definition? Also, are there any symmetric monoidal category where it doesn't commute with coproduct?

\begin{exmp}
    Let $\mathcal{M}(j):=\kappa[\Sigma_j]:= \coprod_{\psi\in\Sigma_j}\kappa_\psi$, coproduct of $\#\Sigma_j$ copies of $\kappa$ (each $\kappa_\psi$ is the unit $\kappa$, subscript is for labeling convenience). It has a natural right $\Sigma_j$-action, where for each $\sigma\in\Sigma_j$, the identity $\kappa_{\psi\sigma^{-1}}\xrightarrow{\sim} \kappa_{\psi}$ induces an isomorphism on coproduct $\tilde{\sigma}\in \Aut(\mathcal{M}(j))$:
\[\xymatrix{
    \kappa_{\psi} \ar[d]_-{\iota_{\psi}} \ar[r]^-{\sim} & \kappa_{\psi\sigma} \ar[d]^-{\iota_{\psi\sigma}}\\
    \coprod_{\psi\in\Sigma_j}\kappa_{\psi} \ar@{-->}[r]_-{\exists ! \tilde{\sigma}} & \coprod_{\psi\in\Sigma_j}\kappa_{\psi\sigma} 
}\]
This construction enforces $\tilde{(\sigma\sigma')} = \tilde{\sigma'}\circ \tilde{\sigma}$, which is a group homomorphism $\Sigma_j^\Op\rightarrow\Aut_\mathcal{V}(\mathcal{M}(j))$.

There is a natural unit morphism $\eta:\kappa\rightarrow \kappa[\Sigma_1] = \kappa$ via identity.

Also, each identity morphism $\kappa_\psi\rightarrow \kappa$ induces a trivialization $t:\kappa[\Sigma_j]\rightarrow\kappa$. Combining the action $\psi\in \Sigma_k$ (by identifying concatination of blocks $(j_{\psi^{-1}(1)},...,j_{\psi^{-1}(k)})=\{1,...,j\}$ in order, also inducing inclusion $\Sigma_{j_s}\hookrightarrow \Sigma_j$), we get the following morphism:
\[\kappa_\psi \tensor \left(\bigtensor_{s=1}^k \kappa[\Sigma_{j_s}]\right)\xrightarrow{\sim} \bigtensor_{s=1}^k \kappa[\Sigma_{j_{\psi^{-1}(s)}}] \xrightarrow{t^{k-1}} \kappa[\Sigma_{j_{\psi^{-1}(s)}}] \rightarrow \kappa[\Sigma_j]\]
Running through all $\psi\in\Sigma_k$, this induces a morphism $\gamma:\coprod_{\psi\in\Sigma_k}\left[\kappa_\psi \tensor \left(\bigtensor_{s=1}^k \kappa[\Sigma_{j_s}]\right)\right] \rightarrow \kappa[\Sigma_j]$, and use the fact $\tensor$ commutes with colimits, this is equivalent to a morphism $\gamma:\kappa[\Sigma_k]\tensor \left(\bigtensor_{s=1}^k \kappa[\Sigma_{j_s}]\right)\rightarrow \kappa[\Sigma_j]$.
\end{exmp}

This is helpful when defining an alternative definition.

Use $\Sigma$ to denotes the category of finite sets, each $n\in\NN$ has $\mathbf{n}:=\{1,...,n\}$ (and $\bzero$ denotes the emptyset), with homsets being all possible isomorphisms (so $n\neq m$ there's no morphism between them, while the endomorphism is consists of permutations). Each object $\mathbf{n}$ has a canonical $\Sigma_n$-action. This is called a \emph{permutation category}.

\hfil

\begin{defn}
    We define a \emph{$\Sigma$-object} in the monoidal category $\mathcal{V}$ to be a functor $F:\Sigma^\Op\rightarrow \mathcal{V}$, which each $F(n)$ is equipped with a natural right $\Sigma_n$-action.

    A morphism between $\Sigma$-objects will be natural transformations, which by definition each morphism at $\mathbf{n}$ is a $\Sigma_n$-equivariant morphism, as all $\sigma\in\Sigma_n$ (as a right action, equivalently an isomorphism) satisfies the following condition:
    \[\xymatrix{
        F(\mathbf{n}) \ar[d]_-{\mu_n} \ar[r]^-{\sigma} & F(\mathbf{n}) \ar[d]^-{\mu_n}\\
        G(\mathbf{n}) \ar[r]_-{\sigma} & G(\mathbf{n})
    }\]
    this forms a \emph{category of $\Sigma$-objects} for $\mathcal{V}$.
\end{defn}
Now, consider the following operation:
\begin{defn}
    Given $F,G$ two $\Sigma$-objects of $\mathcal{V}$, define the \emph{substitution product} $G\circ F:\Sigma^\Op\rightarrow\mathcal{V}$ as follow:
    \[G\circ F(\mathbf{j}):=\coprod_{j_1+...+j_k=j}G(\mathbf{k})\tensor_{\kappa[\Sigma_k]}\left[\left(\bigtensor_{r=1}^k F(\mathbf{j}_r)\right)\tensor_{\kappa[\Sigma_{j_1\times...\times j_k}]}\kappa[\Sigma_j]\right]\]
    Let's assume we know these are associative due to the coequalizer property :)
\end{defn}
But roughly speaking, an operad is a $\Sigma$-object, that is also a monoid under the above product. (There are more details to work out...)

\hfil

Talking about monads relevant to an operad $\mathcal{C}$ over $\mathcal{V}$, we define it as follow:
\begin{defn}
    Define a functor $C:\mathcal{V}\rightarrow\mathcal{V}$ as follow:
    \[\coprod_{j\geq 0} \mathcal{C}(j)\tensor\kappa[\Sigma_j]\tensor X^j \rightrightarrows\coprod_{j\geq 0}\mathcal{C}(j)\tensor X^j \xrightarrow{\text{coeq}}C(X)\]
\end{defn}
Where given any $\psi\in \Sigma_j$, define the top morphism using $\mathcal{C}(j)\tensor \kappa_\psi\tensor X^j\xrightarrow{\sigma\tensor \id}\mathcal{C}(j)\tensor X^j$ (action on $\mathcal{C}(j)$), which inducesleft action morphism $\mathcal{C}(j)\tensor\kappa[\Sigma_j]\tensor X^j\xrightarrow{\ell}\mathcal{C}(j)\tensor X^j$, then using coproduct universality to induce the morphism.

The bottom morphism is defined using $\mathcal{C}(j)\tensor X^j\xrightarrow{\id\tensor \sigma^{-1}}\mathcal{C}(j)\tensor X^j$ (by sending the $i$th index $X$ to the $\sigma^{-1}(i)$th index $X$), then induce the coproduct morphism using similar logic (and label it as $\mathcal{C}(j)\tensor \kappa[\Sigma_j]\tensor X^j\xrightarrow{r}\mathcal{C}(j)\tensor X^j$). This is formally saying: A right action on $\mathcal{C}(j)$ transfers to an inverse left action on $X^j$.

\hfil

Given a morphism $\psi:X\rightarrow Y$, the morphism $\id\tensor \id \tensor \psi^j:\mathcal{C}(j)\tensor \kappa[\Sigma_j]\tensor X^j\rightarrow \mathcal{C}(j)\tensor \kappa[\Sigma_j]\tensor Y^j$, and $\id \tensor \psi^j:\mathcal{C}(j)\tensor X^j\rightarrow\mathcal{C}(j)\tensor Y^j$ bot satisfies the following:
\[\xymatrix{
    \mathcal{C}(j)\tensor \kappa[\Sigma_j]\tensor X^j \ar[d]_-{\id\tensor\id\tensor \psi^j}\ar[r]^-{\ell,r} & \mathcal{C}(j)\tensor X^j \ar[d]^-{\id\tensor \psi^j}\\
    \mathcal{C}(j)\tensor \kappa[\Sigma_j]\tensor Y^j\ar[r]_-{\ell,r} & \mathcal{C}(j)\tensor Y^j
}\]
As including $\mathcal{C}(j)\tensor \kappa_\sigma\tensor X^j$, the top right route provides the morphism $\mathcal{C}(j)\tensor\kappa_\sigma\tensor X^j\xrightarrow{\sigma\tensor \id\tensor\id}\mathcal{C}(j)\tensor X^j \xrightarrow{\id\tensor \psi^j} \mathcal{C}(j)\tensor Y^j$, equivalent to $\mathcal{C}(j)\tensor X^j\xrightarrow{\sigma\tensor \psi^j}\mathcal{C}(j)\tensor Y^j$, and by chasing through the coproduct property, we realize the following diagram commutes:
\[\xymatrix{
    \mathcal{C}(j)\tensor \kappa_\sigma\tensor X^j \ar[d]_-{\id\tensor\id\tensor\psi^j} \ar[r] & \mathcal{C}(j)\tensor \kappa[\Sigma_j]\tensor X^j\ar[d]^-{\id\tensor\id\tensor\psi^j}\\
    \mathcal{C}(j)\tensor \kappa_\sigma\tensor Y^j \ar[r] & \mathcal{C}(j)\tensor \kappa[\Sigma_j]\tensor Y^j
}\]
Combining the two squares, do some diagram chasing, we realize the original square is commuting (with a pinch of coproduct property on $\kappa[\Sigma_j]$). 

Finally, these induce the morphism between coproducts, and they of course commute with the group action by construction (as permuting $\psi^j$ is equivalent to pre-permuting with $\psi^j$, or compose $\psi^j$ with post-permuting). So, composing with the coequalizing morphism of $Y$, the desired diagram factors through coequalizing morphism of $X$, inducing a morphism $C(X)\rightarrow C(Y)$. Such induced morphism is functorial, hence $C$ is indeed a functor.

\hfil

To show $C$ is a monad, we need to construct two natural transformations:

First, define $\eta':\Id_\mathcal{V}\rightarrow C$ by taking the natural unit morphism $\eta:\kappa\rightarrow\mathcal{C}(1)$, and consider the sequence $\eta'_X:X\cong \kappa\tensor X\rightarrow\mathcal{C}(1)\tensor X\rightarrow\coprod_{j\geq 0}\mathcal{C}(j)\tensor X^j\rightarrow C(X)$. This ``transformation'' is natural, because the morphism after applying $C$ is induced in the same manner.

To define the multiplication transforamtion, there are some slight work to do: Consider applying the construction to the first one, we get the following (quite chaotic) isomorphism:
\[\coprod_{k\geq 0}\mathcal{C}(k)\tensor \left(\coprod_{j\geq 0}\mathcal{C}(j)\tensor X^j\right)^k\cong \coprod_{k\geq 0}\coprod_{j_1+...+j_k=j}\mathcal{C}(k)\tensor\mathcal{C}(j_1)\tensor...\tensor\mathcal{C}(j_k)\tensor X^j\]
This is because the $k$-fold monoidal product of the gigantic coproduct records all possible $k$-combinations of $\mathcal{C}(j_i)$, and the product on the copies of $X$ precisely is $X^{j_1+...+j_k}$. Apply the operadic structure morphism $\gamma$ on each entry of the coproduct, we get the following morphism:
\[\overline{\mu'}_X:\coprod_{k\geq 0}\mathcal{C}(k)\tensor \left(\coprod_{j\geq 0}\mathcal{C}(j)\tensor X^j\right)^k\cong  \coprod_{k\geq 0}\coprod_{j_1+...+j_k=j}\mathcal{C}(k)\tensor \left(\bigtensor_{r=1}^k \mathcal{C}(j_r)\right)\tensor X^j \xrightarrow{\sqcup \gamma \tensor \id} \coprod_{j\geq 0}C(j)\tensor X^j \rightarrow C(X) \]
Which, there are some truely loooong diagram chase to check this factors through $C^2(X)$ (by checking it factors through $\mathcal{C}(k)\tensor (C(X))^2$ on each summand $\mathcal{C}(k)\tensor \left(\coprod_{j\geq 0}\mathcal{C}(j)\tensor X^j\right)^k$, then check again it satisfies the equivariance diagram; there are a lot of nuqnces about when to use the fact $\tensor$ commutes with colimits). This defines $\mu:C^2\rightarrow C$.

\hfil

By skipping some details, let's conclude that $\mu'$ is associative (intuitively, it's deduced from the associativity of tensor product; if the morphism before the factorization satisfies certain property, then passing through coequalizer it shoudl still satisfy it). Also, by definition $\eta'$ does serve as a unit (by the fact $\kappa\tensor C(X)\rightarrow \mathcal{C}(1)\tensor C(X)\rightarrow C^2(X)\rightarrow C(X)$ has passed through the collapsing map $\mathcal{C}(1)\tensor C(X)\rightarrow C(X)$, which cancels the unit, and reform identity).

\hfil

Now, to say the correspondance between the operad and the monad, we compare their categories of algebras:
\begin{thm}
Given an operad $\mathcal{C}$ and its corresponding monad $C$, their categories of algebras are equivalent.
\end{thm}

\begin{proof}
    To do.
\end{proof}


\newpage

\section{Review: 2-Categories (ND, not all details)}

This section follows Emily Riehls' \emph{Categorical Homotopy Theory}.

To properly define notions needed afterward, we need a concrete way of viewing categories enriched over small categories. This more generally follows categories enriched over symmetric monoidal categories, which are categories whose hom functor (or some similarly-behaved functor) are endowed with extra structures relevent to a monoidal category.

\begin{notn}
We'll use $(\mathcal{V},\times,*)$ to denote a symmetric monoidal category here. If monoidally closed, we'll use $\mathcal{V}(\_,\_)$ to denote the internal hom functor. 
\end{notn}

\begin{defn}
    A category $\mathcal{D}$ is \emph{enriched over $\mathcal{V}$}, or called a \emph{$\mathcal{V}$-category}, if it consists of the following data:
    \begin{itemize}
        \item For each pair $x,y\in\mathcal{D}$, there is a \emph{Hom-object} $\mathcal{D}(x,y)\in\mathcal{V}$.
        \item For each $x\in\mathcal{D}$, a \emph{unit morphism} $\id_x:*\rightarrow\mathcal{D}(x,x)$ in $\mathcal{V}$.
        \item For each triple $x,y,z\in\mathcal{D}$, a \emph{composition morphism} $\circ:\mathcal{D}(y,z)\times \mathcal{D}(x,y)\rightarrow \mathcal{D}(x,z)$ in $\mathcal{V}$.
    \end{itemize}
    such that the following diagram commutes when given all $x,y,z,w\in\mathcal{D}$:
    \begin{itemize}
        \item \emph{Associativity of composition}:
        \[\xymatrix{
            \mathcal{D}(z,w)\times\mathcal{D}(y,z)\times\mathcal{D}(x,y) \ar[d]_-{\circ\times 1}\ar[r]^-{1\times \circ} & \mathcal{D}(z,w)\times \mathcal{D}(x,z)\ar[d]^-{\circ}\\
            \mathcal{D}(y,w)\times \mathcal{D}(x,y)\ar[r]_-{\circ} & \mathcal{D}(x,w)
        }\]
        \item \emph{Existence of Two-sided Units:}
        \[\xymatrix{
            \mathcal{D}(x,y)\times \text{*} \ar@{=}[dr] \ar[r]^-{1\times \id_x} & \mathcal{D}(x,y)\times \mathcal{D}(x,x) \ar[d]^-{\circ}\\
            & \mathcal{D}(x,y)
        },\quad\quad \xymatrix{
            \mathcal{D}(y,y)\times \mathcal{D}(x,y)\ar[d]_-{\circ} & \text{*} \times \mathcal{D}(x,y) \ar[l]_-{\id_y\times 1} \ar@{=}[dl]\\
            \mathcal{D}(x,y)
        }\]
    \end{itemize}
\end{defn}

These conditions are mimicing the idea of homsets in a given category. For instance, with this definition by choosing $\mathcal{D}(\_,\_):= \Hom_\mathcal{D}(\_,\_)$, we immediately realize all categories are enriched over $\Sets$. Since more general categories don't have concrete descriptions of objects using sets, we instead need a generalization so that the analogy of ``hom functors'' can land in more general categories. 

One thing that sorts of questions me, is the existence of $\mathcal{D}$ as only an assignment of objects, not defined directly as a functor. I think in general a functor works just fine.

\hfil

\begin{exmp}
    Given a one-object category $\mu$ enriched over $(\Ab,\tensor_\ZZ,\ZZ)$, the unique object $*$ has a unit homomorphism $\id_*:\ZZ\rightarrow m(*,*)$, a multiplication homomorphism $\cdot:m(*,*)\tensor_\ZZ m(*,*)\rightarrow m(*,*)$. 

    Denote $1_m:=\id_* (1)$, then we precisely have: $x\dot 1_m=x$, $1_m\dot x=x$, and $(x+y)\cdot (z+w)=x\cdot z+x\cdot w+y\cdot z+y\cdot w$, using unit and $\ZZ$-bilinearity of the morphism. So, one-object category enriched over $\Ab$ with tensor product as monoidal product, is a unital ring.
\end{exmp}

\begin{exmp}
    Given category of topological space $\Top$, the homset $\Hom_\Top (X,Y)$ is in fact enriched over groupoids: Define morphisms between maps $f,g:X\rightarrow Y$ to be homotopy classes of homotopies, then $\Hom_\Top (X,Y)$ forms a groupoid; the unit is $*\mapsto \id_X\in\Hom_\Top (X,X)$ (with the identity morphism on $*$ maps to the class of null homotopy of $\id_X$), and composition is really the traditional composition.
\end{exmp}

\begin{exmp}
    If consider the category of compactly generaetd spaces $\mathcal{U}$, then the product (i.e. the $k$-ification for original cartesian product) satisfies the following adjunction:
    \[\Hom_\mathcal{U}(X\times K,Y)\cong \Hom_\mathcal{U}(X,Y^K)\]
    where $Y^K:=\Hom_\mathcal{U}(K,Y)$
    Which, it's actually cartesian closed and cocomplete.
\end{exmp}

I remembered in the book it's mentioned that for general topological space it's not cartesian closed, but do need to think about why.

\begin{exmp}
    If $\mathcal{V}$ is monoidally closed + cocomplete, then its internal hom functor $\mathcal{V}(\_,\_)$ induces a structure where $\mathcal{V}$ is self-enriched.

    The unit morphism is identified by the adjunction $\Hom_\mathcal{V}(*\times v,v)\cong \Hom_\mathcal{V}(*,\mathcal{V}(v,v))$, identified as the image of the natural isomorphism $\gamma_v:*\times v\xrightarrow{\sim}v$, denoted as $\gamma_v^\sharp:*\rightarrow\mathcal{V}(v,v)$. 

    Before building up the composition morphism, we need to consider an \emph{evaluation morphism} $\epsilon:\mathcal{V}(u,v)\times u\rightarrow v$. This is constructed using the same adjunction, together with the pullback of identity $1:\mathcal{V}(u,v)\xrightarrow{\sim}\mathcal{V}(u,v)$, or $\epsilon := 1^\flat$. Then, we define the pushforward of the following morphism:
    \[\phi:\mathcal{V}(v,w)\times \mathcal{V}(u,v)\times u\xrightarrow{\id\times \epsilon}\mathcal{V}(v,w)\times v\xrightarrow{\epsilon}w,\quad \circ := \phi^\sharp:\mathcal{V}(v,w)\times\mathcal{V}(u,v)\rightarrow \mathcal{V}(u,w)\]
    Let's temporarily believe this works out :)

    (to do: figure out the details for this one)
\end{exmp}

\hfil

Let $\mathcal{V}$ be a cartesian closed monoidal category that's also bicomplete (so, categorical product $\times$ is our monoidal product, final object $*\in\mathcal{V}$ serves as unit due to the fact $*\times X\cong X$ by checking universality; it has an internal hom functor $\mathcal{V}(\_,\_):\mathcal{V}^\Op\times \mathcal{V}\rightarrow\mathcal{V}$ generating the adjoint to $\times$ when fixing one variable, while all desired limits / colimits exist).

The goal is to generated something similar to the \emph{category of based objects} analogous to the category of based sets / based spaces over $\Sets$ / $\Top$ respectively, with \emph{smash product} as our monoidal product. Let's first check these two categories for an understanding:
\begin{exmp}
    Given $\Sets$ or $\Top$, the category of based spaces is equivalent to the slice category of the final object $*\in\Top$ (the image of $*$ becomes the basepoint). We have each space $(X,x_0),(Y,y_0)$ satisfies wedge product $X\vee Y:= X\sqcup Y\/x_0\sim y_0$, and smash product $X\wedge Y:= X\times Y\/X\vee Y$ using the identification with basepoint. 
    
    Consider the inclusion $X \rightarrow X\times Y$ by $x\mapsto (x,y_0)$, and $Y\rightarrow X\times Y$ by $y\mapsto (x_0,y)$, this generaets a map $X\sqcup Y \rightarrow X\times Y$. Which, we claim the following fibre coproduct diagram is satisfied:
    \[\xymatrix{
        X\sqcup Y \ar[d] \ar[r] & X\times Y \ar[d] \ar@(r,u)[ddr]\\
        \text{*} \ar@(d,l)[drr] \ar[r] & X\wedge Y \ar@{-->}[dr]\\
        && Z
    }\]
    By sending $*$ to the class of $X\vee Y$ in $X\wedge Y$, since the image of $X\sqcup Y$ in $X\times Y$ is precisely $X\vee Y$, the above diagram commutes. Moreover, if maps into $Z$ makes the above diagram commutes, we must have the image of $X\vee Y\subseteq X\times Y$ be the point specified by $*$, hence must uniquely factor through $X\wedge Y$ using the quotient space / set property.
\end{exmp} 

\begin{con}
d
\end{con}

\newpage

\section{Permutative Operads, $\mathcal{F}$-spaces}

Need to know:
\begin{itemize}
    \item 2-categories
    \item work out details for $\mathcal{M}$ and the permutative operad $\mathcal{P}=\mathcal{EM}$ (right now there is some basic idea about how it works).
\end{itemize}

\end{document}